% Livro B — parte, capítulo, subseção e lista de exercícios (prompt 03, §52).
\documentclass[11pt,oneside]{book}
% Preâmbulo comum das fixtures do scan (D48). Fontes Latin Modern com T1, para que a camada de
% texto traga os acentos como caracteres, e babel em português, para que \chapter imprima
% "Capítulo" como um livro brasileiro imprime.
\usepackage[T1]{fontenc}
\usepackage[utf8]{inputenc}
\usepackage[brazil]{babel}
\usepackage{lmodern}
\usepackage{amsmath,amssymb}
\usepackage{tikz}
\usepackage[a4paper,margin=2.5cm]{geometry}
\usepackage{enumitem}
\setlist[enumerate,2]{label=\alph*)}
\makeatletter\@ifundefined{c@chapter}{\newcounter{exemplo}}{\newcounter{exemplo}[chapter]}\makeatother
\newenvironment{exemplo}{\par\medskip\refstepcounter{exemplo}\noindent\textbf{Exemplo \theexemplo.}\ }{\par\medskip}
\newcommand{\blocoexercicios}[1][Exercícios]{\section*{#1}}

\begin{document}
\part{Álgebra}
\chapter{Conjuntos}
\section{Operações}
\subsection{União}
A união de dois conjuntos $A$ e $B$ é o conjunto dos elementos que pertencem a $A$ ou a $B$.
\subsection{Interseção}
A interseção de $A$ e $B$ reúne os elementos comuns aos dois conjuntos.

\blocoexercicios[Lista de exercícios]
\begin{enumerate}[label=\textbf{\arabic*.}]
\item Calcule $\{1,2\}\cup\{2,3\}$.
\item Calcule $\{1,2\}\cap\{2,3\}$.
\item Dê um exemplo de conjuntos disjuntos.
\end{enumerate}

\part{Geometria}
\chapter{Triângulos}
\section{Semelhança}
Dois triângulos são semelhantes quando têm os ângulos correspondentes iguais.

\blocoexercicios[Exercícios propostos]
\begin{enumerate}[label=\textbf{\arabic*.}]
\item Prove que dois triângulos equiláteros são sempre semelhantes.
\item Em um triângulo retângulo de catetos $3$ e $4$, calcule a hipotenusa.
\end{enumerate}
\end{document}
